\section{Related Work}
\label{submission:related_work}

Our work is situated at the intersection of static model merging, test-time adaptation, and dynamic routing for multi-task networks. In this section, we review the existing literature in these areas and contrast them with our learning-theoretic approach.

\subsection{Static Model Merging}
Static model merging aims to combine multiple task-specific expert neural networks into a single multi-task model without any retraining or parameter-expansion costs. 
Early methods, such as Task Arithmetic~\cite{ilharco2022taskvectors}, define task vectors by subtracting the pre-trained base model weights from the fine-tuned task-specific weights, followed by a uniform, linear interpolation. 
To resolve destructive interference and task conflicts during static merging, subsequent works introduced advanced selection heuristics. 
For instance, Fisher Merging~\cite{fishermerging} uses the diagonal Fisher Information Matrix of each expert as a local metric tensor to guide the weighted average. 
RegMean~\cite{regmean} computes closed-form linear projections that minimize activation mean-squared errors. 
TIES-Merging~\cite{yadav2023ties} addresses parameter conflict by resolving sign disagreements and pruning redundant small-magnitude parameter updates before merging. 
While static model merging provides a computationally efficient solution, its offline, uniform parameter configuration cannot adapt to input-dependent shifts, resulting in suboptimal performance on diverse multi-task streams.

\subsection{Dynamic Model Merging \& Test-Time Adaptation}
To address the rigidity of static methods, \emph{dynamic model merging} introduces mechanisms to adjust merging coefficients at test-time. 
A key framework in this paradigm is AdaMerging~\cite{lu2023adamerging}, which dynamically adapts merging coefficients on the fly by minimizing unsupervised prediction entropy on incoming test streams. 
While AdaMerging shows strong performance under homogeneous, slow-shifting test streams, subsequent audits have exposed critical vulnerabilities in online test-time adaptation (TTA). 
Studies such as RegCalMerge and SuiteMerge~\cite{regcalmerge, suitemerge} revealed that unsupervised online entropy minimization is highly susceptible to \emph{transductive overfitting}---wherein the routing coefficients overfit to local stream noise or spurious temporal correlations. 
Furthermore, online adaptation suffers from \emph{sacrificial task bias}, where the optimization collapses to favor easy tasks (such as MNIST) while completely sacrificing more complex or noisy tasks (such as SVHN). 
To combat transductive noise, OFS-Tune (Offline Few-Shot Validation)~\cite{OFS2025} employs a tiny offline validation split to identify static or low-degree polynomial merging trajectories. 
Our work builds upon these insights by providing a formal Rademacher complexity analysis of the dynamic blending hypothesis class, establishing a provable generalization bound that mathematically guarantees robustness against transductive overfitting.

\subsection{Dynamic Routing \& Parameter-Space Blending}
Rather than adapting coefficients globally via online optimization, input-dependent routing networks compute sample-specific blending coefficients during the forward pass. 
The Linear Router and the Layer-wise Low-dimensional Linear Router (L3-Router)~\cite{l3router} use parameter-efficient, input-dependent linear projections to predict merging coefficients for each network layer. 
In contrast, Quantum Wavefunction Superposition Merging (QWS-Merge)~\cite{qwsmerge} models task expert weights as orthogonal eigenstates within a complex Hilbert space, predicting wave-like probability amplitudes using trainable phases and trigonometric functions. 
However, recent systematic deconstructions of QWS-Merge~\cite{pendelton2026deconstructing} have proven that its complex wave-like activation collapses in modern representations (such as CLIP or ViT sandboxes), and that classical, L2-regularized linear routing architectures outperform QWS-Merge by margins exceeding +27\% to +43\%. 
Despite their empirical utility, classical linear routers remain heuristic and do not account for hardware-level constraints. 
In particular, when processing heterogeneous batches at the edge, standard hardware engines must average dynamic coefficients across the batch dimension ($\bar{\alpha} = \frac{1}{B} \sum_{b=1}^B \alpha_b$) to maintain $O(1)$ single-model forward execution. 
This hardware constraint induces \emph{heterogeneity collapse}, causing dynamic linear routers to collapse back to uniform averages and dropping accuracy catastrophically (by up to -13.00\%). 
We address this gap by deriving a task-adaptive covariance-weighted regularizer (CFR) that constrains routing parameters to a smooth, low-dimensional manifold, proving resilient against both transductive overfitting and heterogeneity collapse.
